\section{Introduction} 

For commercial buildings, electricity costs are determined not only by total energy consumption but also by peak power: the demand charge, levied on the single highest power draw of a billing cycle, can account for over 30\% of the total bill \cite{neufeld1987demandcharge,neubauer2014demandcharge}.In addition, utilities operate demand response (DR) programs that provide a recurring capacity incentive to buildings that commit to a Firm Service Level (FSL), a limit on the power draw during grid stress events, and that impose penalties when the cap is violated during a called event~\cite{}. Traditionally, building operators curtail their electricity usage and physically turn off appliances to meet these curtailments~\cite{} which can cause inconvenience to the employees. In this work, we propose another avenue to achieve curtailment by leveraging the employee-owned electric vehicles (EVs) and charging facilities that commercial buildings typically have, incentivizing their use for building load management, while allowing users autonomy to refuse charging at the building. The growing presence of EVs at workplace facilities, along with bidirectional charging, is a resource suited to both objectives: a fleet of behind-the-meter batteries that can charge during low-price periods and discharge during peaks, reshaping the building's net load without curtailing building services~\cite{}. We show that by strategically negotiating with privately owned employee EVs having predictable, regular charging schedules at the building, a building operator can achieve lower monthly peaks (lower demand charges) and increase the FSL curtailment value, thereby increasing DR incentives and further reducing the building's electricity cost. The negotiations are targeted toward trading off a user's flexibility in state of charge (SoC) requirements and departure delays for monetary benefits.

Curtailment using the EVs, however, exposes a structural asymmetry between the building operator and the users. The capacity payment is earned by committing, a day or more in advance, to deliver load reduction during events whose timing the building does not control. Unlike a stationary battery, the EVs that help fulfil this commitment are intermittently available, since EVs arrive and depart according to their users' schedules, and are self-consuming. The energy stored in a vehicle may be expended in subsequent travel, with availability determined by charging decisions made on the preceding days \cite{iccps2026_pv2b}. Moreover, the EVs are owned by users who are self-interested and need to be well compensated for helping the building lower its bill. A building that commits to a lower, more aggressive FSL earns greater capacity revenue but bears the full delivery risk, and it needs to strategize and anticipate the user's behavior.
 
Resolving this asymmetry requires navigating a set of deeply interconnected challenges. First, the commitment is made under forecast uncertainty. The FSL must be committed before building loads, EV arrivals, and, critically, user cooperation for the day in question are realized, so the commitment is made against a forecast distribution rather than a realized schedule.  Second, the system must achieve incentive alignment without perfect user-side knowledge. Users must be compensated for flexibility in departure time and required SoC, and must not be able to profit by misreporting \cite{sen2026negotiations}. This is without any prior information available to the building about a user's preferences or a complete prediction of that user's behavior. Crucially, their negotiated SoC requirements must be met by departure. This leads to the third challenge of charging decisions. The charging actions need to balance the power use in high and low load times of the day. Moreover, the charge in EVs is temporally coupled. The energy buffer available for discharge on a high peak or a DR event day is the outcome of past charging decisions, off-site usage of the EVs, and the recurrence pattern of the EV. \cite{iccps2026_pv2b}. Single-day formulations, widely used in the EV charging literature, do not account for the temporal coupling across days, and the demand charge based on the monthly peak power defines sparse and delayed rewards.
 
Prior work addresses these challenges in isolation \textcolor{red}{(related-work section pending)}. Two-stage DR mechanism design \cite{satchidanandan2022two} formalizes the day-ahead-commitment, real-time-delivery structure, but for resources that are dispatchable and stationary rather than intermittently available, have private information, and are self-consuming. Strategy-proof daily negotiation \cite{sen2026negotiations} prices user flexibility and guarantees incentive compatibility, but it does not have temporal coupling across days, so that stored energy cannot be carried across day boundaries and the discharge depth available on a peak day is limited to the incidental surplus with which users arrive. Real-time V2B charging control with user persistence \cite{iccps2026_pv2b} handles the inter-day EV arrivals and charge planning, but it presumes complete cooperation from EV users and does not provide a structure for incentivizing. The reward is only for the long-horizon, sparse demand charge.  None of these approaches individually enables a building operator to provide a prior FSL commitment while requiring no prior commitment from its users.
 
We close this gap with P-CONSENT, a three-layer architecture with a deliberately asymmetric commitment structure. 
Layer~1, a day-ahead optimization between the building and the grid, determines the committed FSL. The commitment and its associated penalty are borne exclusively by the building. However, as shown later, the system tries to reduce the penalties by incorporating them in user negotiations. 
Layer~2, a negotiation mechanism using Monte-Carlo sampled model predictive control (MC-MPC), generates daily user options consistent with the previously fixed FSL. The layer compensates users for realized peak shaving and FSL penalty avoidance, with battery degradation incorporated into the payment. Users are negotiated with only at each arrival, with no future commitments required. 
Layer~3, another real-time Monte Carlo model predictive controller (MC-MPC) \cite{iccps2026_pv2b}, delivers against the committed FSL and the negotiated user requirements at minimum cost, using future samples and re-optimizing at fixed intervals, using updated building loads and EV status. On off-peak days, the controller over-charges consenting users above their negotiated SoC requirements, and is carried over to the subsequent day \cite{iccps2026_pv2b}. On high-peak or DR event days, it discharges the accumulated excess, along with smart planning for intra-day charge/discharge decisions.
Across the system, information flows unidirectionally, from the FSL commitment to the negotiation mechanism to the charging controller.
Our contributions are: 

\begin{itemize} 

\item \textbf{A formulation for DR participation backed by an intermittently available, self-consuming, temporally coupled resource (EVs) with private information}, with a three-layer decision structure whose within-day information flow is unidirectional (from the FSL commitment to the negotiation mechanism to the charging controller) and whose cross-day loop closes through the persistent SoC buffer.
% (Section~\ref{sec:formulation}). 

% \item \textbf{A characterization of the feasible commitment set}, establishing that the set of FSLs deliverable at a given risk level grows monotonically with the forecast buffer, thereby formalizing the manner in which persistence-aware over-charging expands upstream commitment capacity (Proposition~\ref{prop:monotone}), together with a commitment-depth parameter that prices the building's delivery risk and a newsvendor-type characterization of the risk-neutral optimal commitment, which anchors the depth parameter to a principled operating point. 

\item \textbf{A rigorous user negotiation mechanism} where user offers are weakly strategy-proof, incentive compatible and allow voluntary participation.
% we state precisely which of CONSENT's guarantees transfer unchanged and which require restatement under the over-charge policy (Section~4). 

\item \textbf{An evaluation on real building, tariff, and EV fleet data} from a commercial facility (Nissan Advanced Technology Center, Silicon Valley), demonstrating that the combined architecture weakly dominates building-only, negotiation-only, and persistence-only baselines with respect to both building and user cost, with sensitivity analyses over the cost-sharing parameter $\alpha$ \cite{sen2026negotiations} and the commitment depth (Section~\ref{sec:eval}). 

\end{itemize}